%------------------------------ Ejercicio 7.50 -------------------------------
\setcounter{neje}{49}
\begin{nejer}
Defina la función
$
f \colon \mathbb{R}^2 \to \mathbb{R}
$
mediante
$$
f(x, y) \coloneq \frac{\text{sign}(xy)}{x^2 + y^2}, \quad \text{sign}(z) \coloneq \begin{cases} 
1 & \text{si } z > 0, \\
0 & \text{si } z = 0, \\
-1 & \text{si } z < 0,
\end{cases}
$$
para $(x, y) \neq 0$, y tal que
$
f(0, 0) \coloneq 0
$
Pruebe que $f_x, f_y \colon \mathbb{R} \to \mathbb{R}$ son Lebesgue integrables para todo $x, y \en \mathbb{R}$. Pruebe que las funciones $x \mapsto \int_{\mathbb{R}} f_x \dd{m_1}$ y $y \mapsto \int_{\mathbb{R}} f_y \dd{m_1}$ son Lebesgue integrables y
$$
\int_{\mathbb{R}} \left( \int_{\mathbb{R}} f(x, y) \dd{m_1(x)} \right) \dd{m_1(y)} = \int_{\mathbb{R}} \left( \int_{\mathbb{R}} f(x, y) \dd{m_1(y)} \right) \dd{m_1(x)}
$$
No obstante, pruebe que $f$ no es Lebesgue integrable.
\end{nejer}
\begin{sol}
Primero, para cada $x \en \mathbb{R}$ fijo, consideramos la función parcial $f_x(y) \coloneq f(x,y)$. 
Si $x = 0$, entonces $f_0(y) = 0$ para todo $y$, que es claramente Lebesgue integrable con integral $0$.
Si $x \neq 0$, la definición de $f$ da
\[
f_x(y) = \frac{\text{sign}(x)\,\text{sign}(y)}{x^2+y^2}, \qquad y \neq 0
\]
y $f_x(0)=0$. Notemos que $\text{sign}(xy) = \text{sign}(x)\text{sign}(y)$ para todo $x,y$.
El valor absoluto es $|f_x(y)| = \frac{1}{x^2+y^2}$ para $y \neq 0$, y $0$ en $y=0$. Calculamos su integral:
\[
\int_{\mathbb{R}} |f_x(y)| \dd{m_1(y)}
= \int_{-\infty}^{\infty} \frac{1}{x^2+y^2}\dd{y}
= \frac{\pi}{|x|} < \infty
\]
Por lo tanto, $f_x \en L^1(\mathbb{R})$ para todo $x \en \mathbb{R}$. 
De manera análoga, para cada $y$ fijo, $f_y(x) \coloneq f(x,y)$ es Lebesgue integrable y
$\int_{\mathbb{R}} |f_y(x)| \dd{m_1(x)} = \frac{\pi}{|y|}$.

Ahora, como $f_x$ es una función impar de $y$ (el factor $\text{sign}(y)$ es impar y el denominador $x^2+y^2$ es par), su integral sobre $\mathbb{R}$ es nula:
\[
F(x) \coloneq \int_{\mathbb{R}} f(x,y) \dd{m_1(y)} = \int_{\mathbb{R}} f_x(y) \dd{m_1(y)} = 0, \qquad \forall x \en \mathbb{R}
\]
Similarmente,
\[
G(y) \coloneq \int_{\mathbb{R}} f(x,y) \dd{m_1(x)} = \int_{\mathbb{R}} f_y(x) \dd{m_1(x)} = 0, \qquad \forall y \en \mathbb{R}
\]
Las funciones $F$ y $G$ son idénticamente cero, luego son Lebesgue integrables sobre $\mathbb{R}$ y sus integrales valen $0$. Por consiguiente,
\[
\int_{\mathbb{R}} \left( \int_{\mathbb{R}} f(x,y) \dd{m_1(x)} \right) \dd{m_1(y)}
= 0
= \int_{\mathbb{R}} \left( \int_{\mathbb{R}} f(x,y) \dd{m_1(y)} \right) \dd{m_1(x)}
\]

Finalmente, probamos que $f$ no es Lebesgue integrable en $\mathbb{R}^2$. 
Para casi todo $(x,y) \en \mathbb{R}^2$ (salvo en los ejes, que tienen medida de Lebesgue nula) se tiene $|f(x,y)| = \frac{1}{x^2+y^2}$. 
Usando Fubini en coordenadas polares (Ejercicio 7.47):
\[
\int_{\mathbb{R}^2} |f(x,y)| \dd{m_2(x,y)}
= \int_{0}^{2\pi} \int_{0}^{\infty} \frac{1}{r^2}\, r \dd{r} \dd{\theta}
= 2\pi \int_{0}^{\infty} \frac{1}{r}\dd{r}
= \infty
\]
En consecuencia, $f \noten \mathcal{L}^1(\mathbb{R}^2)$, lo que completa la prueba.
\end{sol}
